\chapter{Standards and Design Constraints}
\label{ch:standards}

This chapter describes the standards that MedLift-3D adheres to (along with the alternatives considered), and examines the economic, environmental, ethical, health and safety, social, political, legal, cultural, and sustainability issues. It concludes by the cost analysis and complex engineering problem mapping.

\section{Compliance with the Standards}

MedLift-3D is research software to be developed as medical-device software, so we adhered to the standards which must be met by such software, although not formal certification for a student project. For every standard we explain how it is applied, the other options that were thought of with their pros and cons, and why the chosen one was selected. The Standards are divided into Software Standards, Hardware and Communication Standards.

\subsection{Software Standards}

The software standards are related to data formats, the software life cycle, risk and quality management, code quality and numerical precision and licensing (\cref{tab:swstd}). The input path reads DICOM directly, and checks its geometry, while intermediate data are stored as compressed NumPy files with a JSON header, which store the volume, masks and metadata of each case in one self-describing file, the format in which these are stored and exchanged in hospitals. We have not been certified to IEC 62304, but we did follow the core practices: each requirement was mapped to a test, all code is version controlled, 69 unit tests are run automatically and any anomaly discovered during software development has been documented (\cref{ch:results}) by implementation of the life cycle processes defined in IEC 62304. The hazard analysis in \cref{tab:risk} is used for the application of ISO 14971 risk management and ISO 13485 quality management is identified as required for a future product. This is intentional: physics, measurement is done in 32-bit precision; only diffusion network is done in bf16 to keep the 32-bit range, avoiding overflow of fp16 in 3D networks. Lastly, all data and libraries have open licence with attribution.

\begin{table}[!htbp]
\centering\small
\caption{Software standards.}
\label{tab:swstd}
\begin{tabularx}{\linewidth}{>{\raggedright\arraybackslash}p{2.9cm} >{\raggedright\arraybackslash}X >{\raggedright\arraybackslash}X >{\raggedright\arraybackslash}X}
\toprule
Standard & How it is applied & Alternatives (pros / cons) & Rationale \\
\midrule
DICOM~\cite{dicom} & Input CT read from DICOM; geometry checked against pylidc (0.000~mm difference) & NIfTI: simpler but loses scanner metadata & Hospitals and PACS store images in DICOM \\
Compressed NumPy with JSON headers & Cases, pairs and results keep their geometry in a header & NIfTI cannot hold per-nodule reader masks with metadata; HDF5 is a heavier dependency & One self-describing file per case \\
IEC 62304~\cite{iec62304} (software life cycle) & Not certified; practices followed: requirements traced to tests (\cref{tab:verify}), version control, unit tests, documented anomalies & Ad-hoc research code: faster but unverifiable & Needed for any clinical use; tracing now costs little \\
ISO 14971~\cite{iso14971} (risk management) & Hazard analysis with mitigations (\cref{tab:risk}) & No risk analysis & Required by regulators \\
ISO 13485~\cite{iso13485} (quality management) & Identified as required for a product; out of scope here & -- & Applies to a device manufacturer \\
PEP 8, pytest & Code style; 69 automated tests & No tests: faster now, bugs later & Readable and checkable team code \\
IEEE 754 & fp32 for physics and measurement; bf16 only inside the network & fp16 overflows in 3D U-Nets; fp64 doubles memory & bf16 keeps the fp32 range \\
Open licences & LIDC-IDRI under CC BY 3.0 with citation~\cite{tcia_policy}; PyTorch (BSD), crepes (BSD-3), Gradio (Apache 2.0) & Proprietary data or libraries: costly, restrictive & Legal reuse with attribution \\
\bottomrule
\end{tabularx}
\end{table}

\subsection{Hardware Standards}

MedLift-3D is software, so the hardware standards concern the computing platform and the X-ray equipment that would produce real data (\cref{tab:hwstd}). All GPU work uses NVIDIA's CUDA platform through PyTorch, because CUDA GPUs are the ones offered by every cloud provider we could afford; AMD's ROCm platform is an open alternative but is rarely available for rent. The safety of the X-ray equipment itself is governed by IEC 60601-1 and its part 2-54 for radiography; these apply to the DTS unit, not to our software, and a clinical deployment would rely on a certified unit. Radiation doses are compared as effective dose following ICRP Publication 103, the standard measure of patient risk.

\begin{table}[!htbp]
\centering\small
\caption{Hardware standards.}
\label{tab:hwstd}
\begin{tabularx}{\linewidth}{>{\raggedright\arraybackslash}p{3.2cm} >{\raggedright\arraybackslash}X >{\raggedright\arraybackslash}X >{\raggedright\arraybackslash}X}
\toprule
Standard & How it is applied & Alternatives & Rationale \\
\midrule
NVIDIA CUDA platform & All GPU computation through PyTorch's CUDA backend & AMD ROCm: fewer cloud offers; CPU only: training would take weeks & Widest availability on rented GPUs \\
IEC 60601-1 / 60601-2-54 (X-ray equipment safety) & Applies to the DTS unit that would produce real data, not to our software & -- & We rely on certified DTS hardware \\
ICRP 103 effective dose~\cite{icrp103} & Doses compared in mSv for CT, low-dose CT and DTS & Air kerma: less meaningful for patients & Standard measure of patient risk \\
\bottomrule
\end{tabularx}
\end{table}

\subsection{Communication Standards}

The communication standards concern how the application is reached and how results are exchanged (\cref{tab:commstd}). The application runs as a web server and is reached over HTTP locally and over an encrypted HTTPS tunnel from outside, so any modern browser can use it without installing software; the usual Gradio share link failed on Kaggle, and SSH port forwarding would require setup on the client. Reports, results and calibration files are written as JSON, which is human-readable and language-neutral, unlike Python pickle files, which are unsafe to load from untrusted sources. Integration with hospital systems would use DICOM networking for images and HL7 FHIR for reports, the standards used by picture archiving and electronic record systems.

\begin{table}[!htbp]
\centering\small
\caption{Communication standards.}
\label{tab:commstd}
\begin{tabularx}{\linewidth}{>{\raggedright\arraybackslash}p{3.2cm} >{\raggedright\arraybackslash}X >{\raggedright\arraybackslash}X >{\raggedright\arraybackslash}X}
\toprule
Standard & How it is applied & Alternatives & Rationale \\
\midrule
HTTP/HTTPS with TLS & Application served locally over HTTP and to the browser through an HTTPS tunnel & Gradio share links (failed on Kaggle); SSH port forwarding (client setup) & Encrypted; works from any browser \\
JSON (RFC 8259) & Reports, results and calibration files & XML: verbose; pickle: unsafe to load & Human-readable, language-neutral \\
DICOM networking, HL7 FHIR & Future integration with hospital PACS and records & Custom APIs: no interoperability & Required in hospital workflows \\
\bottomrule
\end{tabularx}
\end{table}

\section{Design Constraints}

Economics, environmental, ethical, health and safety, social, political, sustainability, legal and cultural restrictions influenced the design. Each will be discussed below in conjunction with design impact and action taken.

\subsection{Economic Constraint}

There was no outside funding for the project, so all calculations had to be within a student budget. Wherever possible, free resources were used (laptop CPU, Kaggle GPUs for pilot and the application), CPU-only steps (download and case building) were delegated to a cheap CPU machine, and GPU training was allocated a fixed compute budget. This, in terms of clinical practice, boils down to economic factors: a DTS capable radiography room is significantly less expensive than a CT scanner, and can examine patients much faster, meaning that follow-up with DTS can reduce the cost of follow-up for both hospitals and patients.

\subsection{Environmental Constraint}

\paragraph{Energy and emissions.} The energy consumption of the L40S GPU was minimal, as was the emissions (assuming the per kWh emissions in the grid are approximately 0.48 kg CO$_2$ globally and approximately 0.6 in Bangladesh), for the entire study was conducted on a single GPU. About 0.004 kWh is consumed for 1 measurement (approx. 45 seconds on the GPU). These are the estimates based on rated power as per~\cite{lannelongue2021green} and not metered values.

\paragraph{Data storage and waste.} Downloaded data (64 GB of DICOM) was deleted after being converted to compact cases and intermediate files which could be regenerated were not backed up. No electronic waste was generated as no hardware was purchased.

\paragraph{Radiation as an environmental factor.} This will decrease the collective dose of radiation for the population because CT follow-up (1.5-6 mSv) is replaced by DTS (0.1 mSv).

\paragraph{Mitigations.} A small model (17.4~M parameters) instead of a large generator; a prism instead of whole-volume generation; fixed compute budgets; CPU work on CPU machines; trained models reused for evaluation and the application instead of retraining; and, for future runs, cloud regions powered by renewable energy.

\subsection{Ethical Constraint}

\paragraph{Professional ethics.} We adhered to IEEE~\cite{ieeeethics} and ACM~\cite{acmethics} codes of ethics including honesty in claiming and results.
\begin{itemize}
  \item \textbf{Honest claims:} all results are obtained from \emph{simulated} DTS, not actual scanners, and \cref{sec:gap} provides a list of the things we don't claim.\item \textbf{Honest results:} hypotheses were given prior to results; and the results in the report were whatever we stumbled upon such as: H1 and H2 were not supported; intervals were so wide as to be useless; detector sensitivity is only 35\% at 2 false positives per scan; failed checks in \cref{ch:results}.\item \textbf{Credit:} all datasets, libraries and ideas we build on are cited, including the data citations mandated by TCIA.\item \textbf{Fairness:} LIDC-IDRI is primarily derived from patients in the United States, and performance has not been determined on other patient populations such as patients with healed tuberculosis scars and is not claimed for such populations.\item \textbf{Data ethics:} Public data used under its licence, no re-identification attempt, Only de-identified data.
\end{itemize}

\paragraph{Professional responsibility.} The team takes responsibility for their choices, and their errors. That we lost a third of the attenuation on the prism in the real data pilot, we found out and corrected ourselves on the lack of depth support in our design on the project before the full run, else we'd have been claiming training losses as successes. The owner of each step (\cref{tab:tasks}) is responsible for the checks for that step. The system is decision support -- diagnosis remains the responsibility of the radiologist.

\subsection{Health and Safety Constraint}

\paragraph{Health impacts.} \emph{Physical:} Lower radiation for each follow-up (DTS $\sim$0.1 mSv, vs $\sim$1.5 mSv for low dose CT or several mSv for diagnostic CT) is important given nodule follow-up is repeated over many years~\cite{rampinelli2017,mascalchi2017}. The most significant physical risk is mismeasurement, which can cause a cancer to grow and not be treated or can cause an unnecessary biopsy. \emph{Mental:} having an indeterminate nodule is stressful, and a reassuring or alarming number aggravates it.

\paragraph{Actionable solutions.} Each number is displayed with a calibrated interval and category probabilities to inform the clinician when a result is uncertain; a flag for low confidence triggers the clinician to consider CT and the results should be explained to the patient as categories instead of numbers. The principle hazards are listed in \cref{tab:risk}.

\begin{table}[!htbp]
\centering\small
\caption{Safety risk analysis (ISO 14971 style).}
\label{tab:risk}
\begin{tabularx}{\linewidth}{>{\raggedright\arraybackslash}p{2.6cm} >{\raggedright\arraybackslash}p{2.9cm} >{\raggedright\arraybackslash}p{2.6cm} >{\raggedright\arraybackslash}X}
\toprule
Hazard & Cause & Effect & Mitigation \\
\midrule
Under\-estimated growth & Prior fills depth wrongly & Delayed treatment & Mass-first reporting; intervals; seen/unseen explanation; CT recommended when flagged \\
Overestimated growth & Vessel or chest wall included & Unnecessary CT or biopsy & Segmentation confined to a plausible extent; leaks reported as not measurable \\
Hallucinated nodule & Prior invents structure & False alarm & Data consistency with X-rays; clinician confirms \\
Missed nodule & Detector sensitivity 35\% at 2 FP/scan & Nodule not suggested & Clinician can click anywhere; detector is an aid \\
Out-of-distribution use & Other population, scanner or dose & Silent errors & Validated setting stated; local validation and recalibration required \\
Over-trust & Automation bias & Clinician stops checking & Intervals and flags shown prominently; user training \\
\bottomrule
\end{tabularx}
\end{table}

\subsection{Social Constraint}

Short term, DTS-based follow-up would help to decrease CT waiting lists and patient travel outside of larger cities. It could enable nodule surveillance and ultimately screening in the longer term wherever CT resource availability is the limiting factor and the most survival benefit is achieved through early diagnosis. This translates to less dose, less cost for the patient, a measurement tool with honest uncertainty for the radiologist and local validation and training, and less cost for the health system for follow-up. Two potential adverse effects are the overusing of the tool, where many more people are followed, and the potential for a divide to develop between those who use the tool in public hospitals and those who don't, which is making their own guidelines; both suggest a need for public-sector implementation and adherence to guidelines.

\subsection{Political Constraint}

Public health policy and procurement are key enablers for health technology adoption. In Bangladesh, imaging equipment and software is mostly bought via the government tender process and large private hospital conglomerates, so the adoption of such equipment would require proof to be accepted by health authorities and incorporation into national cancer-control plans. The deployment of foreign cloud computing for health data is also a policy issue -- data and computation should remain in the country.

\subsection{Sustainability}

\paragraph{Lifecycle plan.} The various steps are: (1) development and verification (this project); (2) retrospective validation on real DTS with CT ground truth from a partner hospital; (3) a prospective reader study; (4) regulatory registration; (5) deployment with monitoring of measurement errors and interval coverage; and (6) periodic recalibration of the conformal layer, which only requires new calibration cases, not retraining, when scanners or populations change.

\paragraph{Three dimensions.} \emph{Environmental:} small model, low per measurement energy, no special hardware. \emph{Economic:} uses one mid-range GPU or low cost cloud instance and is open-source software, no licence fees to be paid. \emph{Social:} demonstrated and validated code for other team to maintain and process that can be extended to other limited angle modalities.

\subsection{Legal Considerations}

\begin{itemize}
  \item \textbf{Data licence:} LIDC-IDRI is distributed by TCIA under the Creative Commons Attribution 3.0 Unported licence, which allows commercial, scientific and educational use with attribution. The TCIA data usage policy forbids any attempt to identify or contact the participants and requires every publication to cite both the collection~\cite{armato2011lidc} and the TCIA repository~\cite{clark2013tcia}~\cite{tcia_policy}. We follow both requirements and made no attempt at re-identification.
  \item \textbf{Medical device regulation in Bangladesh:} medical devices are regulated by the Directorate General of Drug Administration (DGDA). The Drugs and Cosmetics Act 2023~\cite{bddca2023}, passed by Parliament in September 2023, replaced the Drugs Act 1940 and the Drugs (Control) Ordinance 1982 and includes provisions on medical devices. Registration follows the DGDA Registration Guidelines for Medical Devices~\cite{dgda2015}, which are based on the Global Harmonization Task Force guidance and define a medical device to include software intended for diagnosis or monitoring. A clinical version of MedLift-3D would therefore be a registrable medical device; implementing rules under the 2023 Act are still being developed, so the current DGDA requirements must be confirmed at the time of submission.
  \item \textbf{Other markets:} US FDA (510(k) or De Novo), EU Medical Device Regulation~\cite{eumdr} and EU AI Act~\cite{euaiact} (high risk medical AI) would apply, all of which entail risk management~\cite{iso14971}, life-cycle processes~\cite{iec62304} and clinical evidence.\item \textbf{Data protection:} The Personal Data Protection Ordinance 2025~\cite{bdpdpo2025} approved on 9 October 2025 and published in the official gazette on 6 November 2025 provides for citizens as owners of their personal data, specific consent is a requirement for collecting, storing, transferring, or using personal data, health data is classified as sensitive and subject to stricter rules and penalties, and cross-border transfers are restricted; some provisions will come into effect 18 months after publication. For a clinical deployment, you'd need to get consent, comply with these regulations (and GDPR/HIPAA for data collected outside of the country) and ideally compute within the country. Only public data, which has been de-identified, is used in the demonstration.\item \textbf{Liability:} intended for decision support, diagnosis is the responsibility of the clinician.
\end{itemize}

\subsection{Cultural Considerations}

\begin{itemize}
  \item \textbf{Language and literacy:} an interface to be deployed should support Bangla and output results in lay categories, numbers to be made available to specialists.\item \textbf{Family-centred decisions:} Diagnoses are often made with the family, so if there is any uncertainty this should be explained by the clinician not in the form of raw probabilities, but in a way the family can follow.\item \textbf{Stigma:} Lung cancer is associated with smoking and blame -- do not blame and many patients, particularly women, have never smoked and are exposed to cooking smoke.\item \textbf{Local disease patterns:} Healed tuberculosis leaves behind calcified granulomas which look like nodules. Such patients were not included in the training of the model, and therefore, local validation is necessary before using the model, in accordance with the patients' right to a model that is tested on patients similar to themselves.\item \textbf{Trust:} many of the patients trust the doctor more than the machine and it is only fair to keep the clinician in charge.
\end{itemize}

\section{Cost Analysis}

\paragraph{Revenue model (illustrative).} Assume 18 months for the commercial version of validation and registration and target customers to be diagnostic centres and hospitals that have DTS-capable radiography (\cref{tab:revenue,tab:cashflow}).

\begin{table}[!htbp]
\centering\small
\caption{Illustrative revenue (BDT lakh).}
\label{tab:revenue}
\begin{adjustbox}{max width=\linewidth}
\begin{tabular}{lccc}
\toprule
Revenue stream (price) & Year 1 & Year 2 & Year 3 \\
\midrule
Pay per study (BDT 300 per nodule follow-up) & 0 & 18 (\num{6000} studies) & 60 (\num{20000} studies) \\
Hospital subscription (BDT 6 lakh per year) & 0 & 18 (3 hospitals) & 48 (8 hospitals) \\
\midrule
\textbf{Total revenue} & \textbf{0} & \textbf{36} & \textbf{108} \\
\bottomrule
\end{tabular}
\end{adjustbox}
\end{table}

\begin{table}[!htbp]
\centering\small
\caption{Illustrative expenses and cash flow (BDT lakh).}
\label{tab:cashflow}
\begin{adjustbox}{max width=\linewidth}
\begin{tabular}{lccc}
\toprule
Item & Year 1 & Year 2 & Year 3 \\
\midrule
Development team (3 engineers) & 25 & 25 & 30 \\
Clinical validation study & 20 & 10 & -- \\
Regulatory registration and quality system & 8 & 4 & 4 \\
Cloud / GPU and hosting & 3 & 5 & 8 \\
Marketing and training & 2 & 6 & 8 \\
\midrule
\textbf{Total expenses} & \textbf{58} & \textbf{50} & \textbf{50} \\
Net cash flow & $-58$ & $-14$ & $+58$ \\
\textbf{Cumulative cash flow} & $\mathbf{-58}$ & $\mathbf{-72}$ & $\mathbf{-14}$ \\
\bottomrule
\end{tabular}
\end{adjustbox}
\end{table}

Under these assumptions the venture breaks even early in year four. The largest uncertainty is the cost and duration of clinical validation, which is why a partnership with a public hospital is the first step.

\section{Complex Engineering Problem}

MedLift-3D is a complex engineering problem: it is ill-posed by physics, combines X-ray physics, generative AI and statistics, has conflicting requirements (dose, accuracy, honesty, cost), and has direct consequences for patients.

\subsection{Complex Problem Solving}

\begin{table}[!htbp]
\centering\small
\caption{Mapping with complex problem solving.}
\label{tab:cep}
\begin{adjustbox}{max width=\linewidth}
\begin{tabular}{ccccccc}
\toprule
\makecell{P1\\Depth of\\knowledge} & \makecell{P2\\Range of conflicting\\requirements} & \makecell{P3\\Depth of\\analysis} & \makecell{P4\\Familiarity\\of issues} & \makecell{P5\\Extent of\\applicable codes} & \makecell{P6\\Stakeholder\\involvement} & \makecell{P7\\Inter-\\dependence} \\
\midrule
\cmark & \cmark & \cmark & \cmark & \cmark & $\circ$ & \cmark \\
\bottomrule
\end{tabular}
\end{adjustbox}
\end{table}

\paragraph{P1 -- Depth of knowledge.} What is needed extra to the standard curriculum is X-ray projection physics, linear algebra of range and null spaces, probabilistic diffusion models, conformal statistics and 3D deep learning. The knowledge profiles are shown in \cref{tab:wk}.

\begin{table}[!htbp]
\centering\small
\caption{Knowledge profile mapping for P1.}
\label{tab:wk}
\begin{tabularx}{\linewidth}{>{\raggedright\arraybackslash}p{4.2cm} >{\raggedright\arraybackslash}X}
\toprule
Knowledge profile & Where it is used \\
\midrule
WK1 Natural sciences & X-ray attenuation, Poisson photon statistics, radiation dose \\
WK2 Mathematics & Projection moments and conditioning; CGLS; probability and exchangeability for conformal prediction \\
WK3 Engineering fundamentals & Numerical methods (SIRT, conjugate gradients), algorithms, data structures \\
WK4 Specialist knowledge & Tomography, diffusion models, 3D U-Nets, Gaussian representations \\
WK5 Engineering design & Requirements, prism design, verification steps, application \\
WK6 Engineering practice & Cloud GPU operation, testing, version control, reproducibility \\
WK7 Engineering in society & Radiation safety, regulation, ethics, cultural factors \\
WK8 Research literature & Literature on DTS, sparse-view CT, diffusion, Gaussians and conformal prediction \\
\bottomrule
\end{tabularx}
\end{table}

\paragraph{P2 -- Range of conflicting requirements.} The dose vs accuracy, the filling of the missing wedge vs no anatomical invention, the single numbers vs honest intervals, the memory vs the field of view, and the detection recall vs false positives (\cref{tab:conflicts}).

\paragraph{P3 -- Depth of analysis.} We started from the ground up to identify what DTS can measure, created a virtual 583 patient trial and looked at calibration, size group coverage and seen/unseen content not just image quality.

\paragraph{P4 -- Familiarity of issues.} For the calibrated 3D nodule measurement in DTS, no standard solution exists and already the combination of a data-consistent prior in a DTS-prism, Gaussian mass and conformal calibration had to be designed.

\paragraph{P5 -- Extent of applicable codes.} DICOM, IEC 62304, ISO 14971, data licences and medical device regulation (\cref{tab:swstd})

\paragraph{P6 -- Extent of stakeholder involvement (partial).} Clinical needs were based on guidelines and feedback from supervisor; future work will include direct evaluation by the radiologists.

\paragraph{P7 -- Interdependence.} There is a lot of interdependence between simulation, reconstruction, prior, Gaussians, calibration, and detector: for instance, until now the depth-support bug in reconstruction screwed up the training data, and consequently all subsequent measurements, including calibration and detector.

\subsection{Engineering Activities}

\begin{table}[!htbp]
\centering\small
\caption{Mapping with complex engineering activities.}
\label{tab:cea}
\begin{adjustbox}{max width=\linewidth}
\begin{tabular}{ccccc}
\toprule
\makecell{A1\\Range of resources} & \makecell{A2\\Level of interaction} & \makecell{A3\\Innovation} & \makecell{A4\\Consequences for\\society and environment} & \makecell{A5\\Familiarity} \\
\midrule
\cmark & \cmark & \cmark & \cmark & \cmark \\
\bottomrule
\end{tabular}
\end{adjustbox}
\end{table}

\paragraph{A1 -- Range of resources.} There was a public collection of 583 patients ($\approx$64 Gb), rented CPU and GPU machines with shared storage, free Kaggle GPUs and open-source libraries.

\paragraph{A2 -- Level of interaction.} Three persons in the team (with different roles), the supervisor, and the course teacher, published clinical guidelines, and open-source communities (crepes, Gradio, pylidc) interact.

\paragraph{A3 -- Innovation.} The mass-first projection-moment argument for DTS, the prism-based data-consistent diffusion prior, decision-level calibrated uncertainty and a seen/unseen test with respect to the truth.

\paragraph{A4 -- Consequences for society and environment.} The intervals, flags and clinician control all address the potential for incorrect measurements, while providing lower radiation and cost per follow-up and greater access to nodule surveillance (\cref{tab:risk}).

\paragraph{A5 -- Familiarity.} During the project, diffusion models, conformal predictive systems and null-space analysis were learned, which are much more than typical undergraduate projects.

\section{Summary}

MedLift-3D follows the software, hardware and communication standards relevant to medical imaging software and prepares for IEC 62304 and ISO 14971. Its main impacts are positive (lower dose and cost, wider access), with clear risks (wrong measurements, over-trust, untested populations) that the design mitigates through calibrated intervals, flags, clinician control and a validation plan. It meets the attributes of a complex engineering problem.
